\section{Boolean growth at a fixed real displacement}\label{sec:boolean}

Fix a prime drop from lag \(ph\) to lag \(h\), and write \(H=3^h\),
\(M=[p]_H=1+H+\cdots+H^{p-1}\). We work with actual points
\(x\in C\cap(0,1)\) of reduced denominator \(q\) prime to \(3\);
the two integer endpoints can be counted separately. Fix the signed
\emph{real} number \(W=\{H^p x\}-x\), where braces denote fractional
part. Further restrictions on the possible witnesses are allowed.
Write \(d\) for the denominator of \(W\) (with denominator \(1\) for
zero). For each integer \(u\) let \(A_u\) consist of the distinct
circle values \(w=(H-1)x\bmod 1\) of exact denominator \(ud\).
In fact the denominator of each such \(w\) is a multiple of \(d\):
\(Mw=W\bmod 1\), so \(u=\gcd(M,ud)\). In particular \(u\mid M\);
two such values differ in the kernel of multiplication by \(M\) on
\(G_{ud}=(ud)^{-1}\mathbb Z/\mathbb Z\), which is \(G_u\). Thus
\(A_u\) is contained in one affine \(G_u\)-grid. The real child
\(\{Hx\}-x\in(-1,1)\) has at most two values above any given \(w\).

\begin{proposition}[Growth of a conditional winding]\label{prop:bool-winding}
For every nonempty \(A_u\) as above,
\[
 |A_u|^{1/\kappa}
   \ll p^{3/2}u^{1/(p-1)}\,|A_u-A_u|,
 \qquad \kappa=\log 2/\log 3.
\]
The implied constant is absolute, independent of \(h\) and of the
restrictions on the witnesses.
\end{proposition}

\begin{proof}
Select one witness \(x=x(w)\) for each \(w\). Set
\[
 m=1+\Big\lceil \log_3(4\sqrt p)+\frac{\log_3 u}{p-1}\Big\rceil,
 \qquad x_j=\{H^j x\}\quad(0\le j\le p).
\]
The points \(x_j\) have their unique ternary expansions in the
digits \(0,2\): multiplication by \(H\) removes \(h\) digits, and
their reduced denominators remain prime to \(3\) and exceed one.
Form a Boolean vector \(F(w)\) by recording the first \(m\) ternary
digits, divided by \(2\), at each \(x_j\), \(0\le j<p\).

We first prove injectivity. Suppose \(F(w)=F(w')\), and use primes
on \(x_j\) for the second witness. Put \(v_j=x_j-x'_j\); the prefix
bound gives \(|v_j|\le 3^{-m}\) for \(j<p\). The equality
\(x_p-x_0=x'_p-x'_0=W\) gives \(v_p=v_0\) as real numbers.
Consequently \(t=(v_{j+1}-v_j)_{j=0}^{p-1}\) lies in the real
hyperplane \(E=\{y:\sum_j y_j=0\}\) and has Euclidean norm
\(\le 2\sqrt p\,3^{-m}\). Modulo \(\mathbb Z^p\), its coordinates
are \((H^j(w-w'))_{0\le j<p}\).

Consider the lattice
\[
 \Lambda=\{y\in E:y\bmod\mathbb Z^p
             =(H^j z)_{0\le j<p}\text{ for some }z\in G_u\}.
\]
It contains \(E\cap\mathbb Z^p\) with index at most \(u\), by the
first coordinate of the reduction map. The latter integer lattice
has covolume \(\sqrt p\) in \(E\) (use the basis \(e_j-e_{p-1}\),
\(0\le j<p-1\)). Thus \(\operatorname{covol}(\Lambda)\ge 1/u\).
Since \(u\mid M\), \(H^p z=z\) on \(G_u\); cyclic coordinate shift
preserves \(\Lambda\). If \(w\ne w'\), then \(t\ne0\) is a rational
vector in \(\Lambda\). The shift \(R\) on \(E\cap\mathbb Q^p\)
satisfies \(1+R+\cdots+R^{p-1}=0\). As
\(\Phi_p(X)=1+\cdots+X^{p-1}\) is irreducible over \(\mathbb Q\),
the minimal annihilating polynomial of this nonzero \(t\) is
\(\Phi_p\). Hence \(t,Rt,\ldots,R^{p-2}t\) span a full-rank
sublattice, of covolume at least \(1/u\). Hadamard's inequality
instead bounds their volume by \((2\sqrt p\,3^{-m})^{p-1}<1/u\).
This proves injectivity. Fixing the prior only modulo one would
not give the real equality \(v_p=v_0\) used here.

Next fix a single circle difference \(z=w-w'\). Put
\(e_j=(x_j-x'_j)/2\in(-1/2,1/2)\).
The integers removed when passing from \(H^j x\) to \(x_j\)
have only even ternary digits and are even; the same holds for
\(x'\). Thus \(e_j\equiv H^j e_0\pmod1\).
Writing \(a=(H-1)e_0\bmod1\), we have \(2a=z\), with at most two
choices of \(a\), and
\[
 e_j\equiv e_0+[j]_H a\pmod1,\qquad [0]_H=0.
\]
The ternary digits of \(1/2+e_j\) are one plus the successive
differences of the bits of \(x_j,x'_j\), by the series
\(1/2=\sum_{\ell\ge1}3^{-\ell}\). This number belongs to \((0,1)\)
and has denominator prime to \(3\), so it is not a ternary interval
endpoint. For a fixed \(a\), as the single parameter \(e_0\bmod1\)
varies, the vector of length-\(m\) prefixes of all \(1/2+e_j\) is
constant between at most \(p3^m\) endpoints: these are translates
of the \(3^m\) endpoints of the ternary partition, one partition for
each \(j\). No witness lies on them. The total number of differences
\(F(w)-F(w')\) over this \(z\) is therefore at most \(2p3^m\).

For \(E_1,E_2\subset\{0,1\}^s\) the integer-cube sumset inequality
\cite[Corollary~1.11]{BIMP} gives
\(|E_1+E_2|\ge (|E_1||E_2|)^{(\log_2 3)/2}\). Apply it to
\(F(A_u)\) and its translate-reflection \(\mathbf1-F(A_u)\). Then
\[
 |A_u|^{1/\kappa}\le |F(A_u)-F(A_u)|
     \le 2p3^m |A_u-A_u|.
\]
The definition of \(m\) gives the asserted coefficient.
\end{proof}

\paragraph{The terminal fiber.}
At the final lag \(h\), fix \(W=\{3^h x\}-x\) of denominator \(d\).
For a fixed \(u=\gcd(q,3^h-1)\), the relation
\((3^h-1)x=W\bmod1\) gives \(q=ud\). Its solutions in \(G_q\)
lie in one coset of \(G_u\). Now let \(A_u\) be the set of genuine
points \(x\) on this fiber, read modulo one. We have the stronger
bound
\begin{equation}\label{eq:bool-terminal}
 |A_u|^{1/\kappa}\le 2|A_u-A_u|.
\end{equation}
Indeed take \(m\) with \(3^{-m}<1/u\). Distinct points of one affine
\(G_u\)-grid cannot share their first \(m\) digits. For each fixed
circle difference there are at most two real differences \(x-x'\)
in \((-1,1)\); each specifies one Boolean prefix difference by the
ternary expansion of \(1/2+(x-x')/2\), as above. The same cube
inequality proves \eqref{eq:bool-terminal}.

\section{Stabilization and the quotient estimate}\label{sec:lcm}

Work at one fixed parent, either a drop or a terminal fiber as
above, with fixed common period \(n=\operatorname{ord}_q(3)\)
and one bin \(U/2<u\le U\). The reduced denominators \(q\le T\)
of different witnesses need not coincide. In both cases \(u\mid q\),
so \(\operatorname{ord}_u(3)\mid n\) (take the period to be \(1\)
if \(u=1\)); the represented-period estimate gives \(n=T^{o(1)}\)
uniformly. Discard empty \(A_u\). We distinguish \(mS\), the sum
of \(m\) copies of a set, from \(m\cdot S\), its image under scalar
multiplication. Define the symmetric, exactly tripling-invariant set
\[
 S_u=\bigcup_{0\le t<\operatorname{ord}_u(3)}3^t\cdot(A_u-A_u)
       \subset G_u .
\]
It contains zero. No invariance of \(A_u\) is claimed. For every
\(\theta>0\) choose an integer \(a>1/\theta\). For some
\(j\in\{1,10,\ldots,10^{a-1}\}\) we have
\(|10jS_u|\le U^\theta |jS_u|\): otherwise iteration would give
\(|10^a S_u|>U\). Put \(J=10^{a-1}\), and let \(B_u=4jS_u\) for
one such \(j\). The set \(2B_u=8jS_u\) is covered by at most
\(U^\theta\) translates of \(B_u\). Indeed, take a maximal family
of disjoint translates \(y+jS_u,\ y\in 8jS_u\); all lie in
\(9jS_u\subset10jS_u\). Each unchosen translate meets a chosen
one, so each \(y\) is in a chosen center plus \(2jS_u\subset B_u\).
The set \(B_u\) remains symmetric and exactly tripling-invariant,
and contains \(A_u-A_u\). We obtain, at a drop with \(p>P_0\),
\begin{equation}\label{eq:bool-container}
 |A_u|\ll T^{o(1)} U^{\kappa/(P_0-1)} |B_u|^\kappa ,
\end{equation}
using \(p\le n\); at the terminal the factor preceding
\(|B_u|^\kappa\) can be replaced by \(2^\kappa\).

\begin{lemma}[Quotient inventory at fixed complement]\label{lem:lcm}
Fix \(\theta,\sigma'>0\) and \(\eta>0\), and let
\(T^\eta\le U\le 2T\). For each integer \(1\le K\le U\), put
\(V=U/K\). At either kind of parent, for any subcollection of its
allowed indices,
\[
 \sum_{D:\,KD\ {\rm in\ the\ subcollection}} |K\cdot B_{KD}|
     \ll V U^{\sigma'} T^{o(1)} .
\]
For a drop the bound requires \(p>P_0(\theta,\sigma')\);
there is no time-prime condition at the terminal. The error term
is uniform in the parent, \(K\), the subcollection and the chosen
members of the finite menu. No coprimality of \(K,D\) is assumed.
\end{lemma}

\begin{proof}
Choose a fixed integer \(k>2/\sigma'\). For an ordered tuple of
admissible \(D_i\le V\), write
\[
 L=\operatorname{lcm}(D_1,\ldots,D_k),\qquad
 Q=K L=\operatorname{lcm}(KD_1,\ldots,KD_k).
\]
The identity follows prime by prime, including at primes dividing
\(K\). In the drop case \(Q\mid[p]_H\), for the \emph{same}
\(H,p\) throughout this tuple. Each constituent of
\(\sum_i B_{KD_i}\) can be expanded using the defining union of
rotations. For a choice of rotations the result is a sum of at
most \(4Jk\) sets of the form \(X-X\), where \(X=3^t\cdot A_{KD_i}\)
lies on one affine \(G_Q\)-grid. Indeed \(G_{KD_i}\subset G_Q\);
also \(3^t(H-1)x=(H-1)\{3^t x\}\bmod1\) with
\(\{3^t x\}\in C\). At the terminal \(X\) consists directly of
genuine shifted points \(\{3^t x\}\in C\). Thus all the difference
sets are of the kind tested by the common mixed-fiber estimate.
Their grid offsets may vary independently; the copies need not
share a realized prior. Since taking differences removes every
individual recentering, there is no cost for aligning offsets.

There are at most \(n^{4Jk}\) choices of the rotations. Let \(S\)
be a dyadic detector height with \(Q\le S\). For any fixed
\(\delta>0\), Theorem~\ref{thm:native},
at the finitely many orders up to \(4Jk\), gives
\begin{equation}\label{eq:lift-lcm}
 \left|\sum_{i=1}^k K\cdot B_{KD_i}\right|
 \le \left|\sum_{i=1}^k B_{KD_i}\right|
 \ll n^{4Jk} S^\delta D_S(Q).
\end{equation}
Here \(D_S(Q)\) is the common full-grid detector of
Proposition~\ref{prop:detector}, whose sum over
\(Q\le S\) is \(\le S^{1+o(1)}\). The first inequality is the
cardinality bound for the image under the \emph{single} scalar
\(K\); it uses no division by \(H-1\). At the terminal we use the
true-\(C\) assertion of the same theorem.

On the output side \(K\cdot B_{KD_i}\subset G_{D_i}\).
Addition from \(\prod_i G_{D_i}\) onto \(G_L\) is surjective:
its image is generated by the \(1/D_i\), and
\(\gcd_i(L/D_i)=1\). Each fiber has exactly \(\prod_i D_i/L\)
points. Restricting to the indicated subsets and applying
\eqref{eq:lift-lcm} gives
\[
 \prod_i |K\cdot B_{KD_i}|
       \ll \frac{V^k}{L}\,n^{4Jk} S^\delta D_S(KL).
\]
There are at most \(\tau(L)^k=T^{o(1)}\) ordered tuples for each
\(L\), by the divisor bound since \(L\le V^k\le(2T)^k\).
Group by \(W\le L<2W\) for dyadic \(W\ge1\), taking a \emph{single}
dyadic \(S\) with \(2KW\le S<4KW\) for that shell. Nonempty shells
have \(S\ge Q>U/2\); also \(S\le 4 K V^k\le 4U^k\). Since \(KL\)
is different for each \(L\), the complete detector budget, even
though it includes many unused moduli, gives
\[
 \sum_{W\le L<2W}D_S(KL)\le S^{1+o(1)}.
\]
After expansion of the \(k\)-th power of the desired sum, the shell
therefore contributes at most
\[
 V^k T^{o(1)}\,\frac{S^{1+\delta}}{W}
   \ll_{k,\delta} V^k K U^{k\delta} T^{o(1)} .
\]
The subpower factor includes \(n^{4Jk}\) and the uniform detector
budget error. There are \(O_k(\log(2T))\) shells. Taking a \(k\)-th
root costs \(K^{1/k}\le U^{1/k}\), not a factor \(K\) in each
coordinate. Choose \(\delta<\sigma'/2\), then \(P_0\) sufficiently
large for the fixed mixed orders and this tolerance, and finally
\(T\) large. This proves the lemma, including \(K=1\) and
\(D_i=1\). It prices the \(D\)'s for one specified \(K\); it does
not sum over possible complements.
\end{proof}
